\section{April 13}


\subsection{The Kronecker-Weber Theorem}

  \begin{proof}[Proof of \Yang{Kronecker-Weber Theorem}]
    Pick $\sigma \in \Gal(K^\ab/K)$ such that $\sigma|_{K_{K_\pi K^\unr}} = \phi_\pi(\pi)$ and $F := (K^\ab)^\sigma$. 
    As $\sigma|_{K_\pi} = \phi_\pi(\pi)|_{K_\pi} = \id_{K_\pi}$ and $\sigma|_{K^\unr} = \phi_\pi(\pi)|_{K^\unr} = \Frob_K$,
    we have $K_\pi \subset F$ and $F \cap K^\unr = K$.
    By \Yang{Lemma}, we get $F = K_\pi$.
    Hence we only need to show $F \cdot K^\unr = K^\ab$.

    For all $M/F$ finite extension in $K^\ab/F$, consider
    \[ \tau_M : \hat{\bbZ} \to \Gal(M/F), \quad n \mapsto \sigma^n|_M. \]
    It is surjective by definition of $F$.
    Take inverse limit of $M$, we get 
    \[ \tau: \hat{\bbZ} \surjmap \invlim_{M} \Gal(M/F) = \Gal(K^\ab/F). \]
    \Yang{need compactness of $\hat{\bbZ}$.}
    Consider the following commutative diagram:
    \[ \begin{tikzcd}
        \hat{\bbZ} \ar[r, "\tau", two heads] \ar[rd, "n \mapsto \Frob_K^n"', hook] & \Gal(K^\ab/F) \ar[d, "g \mapsto g|_{K^\unr}"] \\
        & \Gal(K^\unr/K).
    \end{tikzcd} \]
    Hence the restriction map $\Gal(K^\ab/F) \to \Gal(K^\unr/K)$ is injective.
    That is, $\Gal(K^\ab/F) \cap \Gal(K^\ab/K^\unr) = \{1\}$.
    Hence $\Gal(K^\ab/(F \cdot K^\unr)) = \{1\}$, which implies $F \cdot K^\unr = K^\ab$.
  \end{proof}

  \begin{corollary}
    The maximal abelian extension of $\bbQ_p$ is obtained by adjoining all roots of unity.
  \end{corollary}

  \begin{example}
    A finite abelian extension need not be of form $L_t L_u$ with $L_t/K$ totally ramified and $L_u/K$ unramified.
     
    For example, take $\bbQ_5[\xi_5,\xi_{5^4-1}]/\bbQ_5$.
    Note $\bbQ_5[\xi_5]/\bbQ_5$ is totally ramified and $\bbQ_5[\xi_{5^4-1}]/\bbQ_5$ is unramified, 
    and $\Gal(\bbQ_5[\xi_5]/\bbQ_5) \cong (\bbZ/5\bbZ)^\times = \left<\tau\right>$ is of order $4$ while $\Gal(\bbQ_5[\xi_{5^4-1}]/\bbQ_5) \cong \left<\sigma\right>$ with $\sigma^4 = 1$.
    Then $\Gal(\bbQ_5[\xi_5,\xi_{5^4-1}]/\bbQ_5) \cong \left<\tau\right> \times \left<\sigma\right>$, where 
    \[ \begin{cases}
        \tau(\xi_5) = \xi_5^2, \quad \tau(\xi_{5^4-1}) = \xi_{5^4-1}, \\
        \sigma(\xi_5) = \xi_5, \quad \sigma(\xi_{5^4-1}) = \xi_{5^4-1}^5.
    \end{cases} \]
    Let $L$ be the fixed field of $\sigma^2 \tau$.
    Then $[L:\bbQ_5] = 4$.
    Let $L_u$ be the maximal unramified subextension of $L/\bbQ_5$, i.e. $L_u = L \cap \bbQ_5[\xi_{5^4-1}] = \bbQ_5[\xi_{5^2-1}]^{\sigma^2 \tau} = \bbQ_5[\xi_{5^2-1}^{25}]$.
    Hence $L_u/\bbQ_5$ is of degree $2$.
    If $L = L_t L_u$ with $L_t/\bbQ_5$ totally ramified for some $L_t$, 
    then $\Gal(L/\bbQ_5) \cong \Gal(L_t/\bbQ_5) \times \Gal(L_u/\bbQ_5)$.
    However, $\Gal(L/\bbQ_5) \cong \bbZ/4\bbZ$ is cyclic, which is a contradiction.
  \end{example}

  \begin{theorem}[Kronecker-Weber Theorem]
    Every abelian extension of $\bbQ$ is contained in a cyclotomic extension.
  \end{theorem}
  \begin{proof}
    Let $K/\bbQ$ be a finite abelian extension, $p$ a prime number, $\frakp$ a prime of $K$ above $p$ which is ramified over $\bbQ$.
    
    The \Yang{Local Kronecker-Weber Theorem} implies $K_\frakp \subset \bbQ_p[u_{p,\frakp},v_{p,\frakp}]$ for some roots of unity $u_{p,\frakp},v_{p,\frakp}$ with $u_{p,\frakp}$ of $p$-power order and $v_{p,\frakp}$ of order prime to $p$.
    Let $L = \bbQ[{u_{p,\frakp}}]$ for all such $(p,\frakp)$.
    Let $K' = K \cdot L$.
    Then $K'/L$ is still abelian.
    For all prime $\frakP$ of $K'$ over $\frakp$, we have $K'_\frakP \subset \bbQ_p[u_{p,\frakp}, w_{p,\frakP}]$ for some root of unity $w_{p,\frakP}$ of order prime to $p$.
    We may replace $K$ by $K'$ and assume $L \subset K$.

    We have 
    \[ [K:\bbQ] \geq [L:\bbQ] = \prod_{p \text{ ramified}}\phi(p^{s_p}). \]
    Since $G$ is abelian, the inertia subgroup $I_\frakp$ depends only on $p$.
    On the other hand, $K_\frakp \subset \bbQ_p[u_{p,\frakp}, v_{p,\frakp}]$ implies $\#I_p \leq \phi(p^{s_p})$.
    By \cref{lem:inertia-generate-Galois-over-Q}, we have $G$ is generated by the inertia subgroups $I_p$ for all ramified primes $p$.
    Hence 
    \[ \# G \leq \prod_{p \text{ ramified}} \#I_p \leq \prod_{p \text{ ramified}} \phi(p^{s_p}) \leq [L:\bbQ] \leq [K:\bbQ] = \# G. \]
    It enforces all inequalities to be equalities, which implies $[K:\bbQ] = [L:\bbQ]$ and hence $K = L$ is cyclotomic.
  \end{proof}

  \begin{lemma}\label{lem:inertia-generate-Galois-over-Q}
    Let $K/\bbQ$ be a finite Galois extension with $G = \Gal(K/\bbQ)$.
    Then $G$ is generated by the inertia subgroups $I_\frakp$ for all primes $\frakp$ of $K$ that are ramified over $\bbQ$.
  \end{lemma}
  \begin{proof}
    Let $H$ be the subgroup of $G$ generated by all inertia subgroups $I_\frakp$ for all primes $\frakp$ of $K$ that are ramified over $\bbQ$.
    Let $M = K^H$ be the fixed field of $H$.
    Then for all prime $\frakp$ of $K$, $K^I_\frakp \supset M$.
    Then $\frakp \cap M$ is unramified over $\bbQ$.
    Hence $M/\bbQ$ is unramified at all primes.
    However, the only unramified extension of $\bbQ$ is $\bbQ$ itself.
    Hence $M = \bbQ$, which implies $H = G$.
  \end{proof}



\subsection{Cohomology of Groups}

  Let $G$ be a group.

  \begin{definition}
    A \emph{$G$-module} is a $\bbZ[G]$-module, i.e., an abelian group $M$ equipped with an action of $G$ that is compatible with the group structure.
    A \emph{morphism} of $G$-modules is a $\bbZ[G]$-module homomorphism, i.e., a group homomorphism that commutes with the $G$-action.
    We denote the category of $G$-modules by $\Mod_G$.
  \end{definition}

  Let $M,N$ be $G$-modules.
  Then $\Hom_{\Ab}(M,N)$ is also a $G$-module with the action defined by 
  \[ (g \cdot \varphi)(m) := g \cdot \varphi(g^{-1} \cdot m). \]

  \begin{definition}
    Let $H < G$ be a subgroup and $M$ a $H$-module.
    The \emph{induced module} $\upoperatorname{Ind}_H^G M$ is defined as
    \[ \upoperatorname{Ind}_H^G M := \{\phi: G \to M \mid \phi(hg) = h \cdot \phi(g) \text{ for all } h \in H, g \in G\}. \]
    The $G$-action on $\upoperatorname{Ind}_H^G M$ is given by 
    \[ (g \cdot \phi)(x) := \phi(xg) \text{ for all } g,x \in G. \]
    
    A homomorphism $\alpha: M \to M'$ of $H$-modules induces a homomorphism $\upoperatorname{Ind}_H^G \alpha: \upoperatorname{Ind}_H^G M \to \upoperatorname{Ind}_H^G M'$ of $G$-modules defined by 
    \[ (\upoperatorname{Ind}_H^G \alpha)(\phi) := \alpha \circ \phi \text{ for all } \phi \in \upoperatorname{Ind}_H^G M. \]
  \end{definition}

  \begin{remark}
    We have a natural isomorphism 
    \[ \upoperatorname{Ind}_H^G M \cong \prod_{[\alpha] \in H \backslash G} \alpha^{-1} M, \]
    with the $G$-action given by 
    \[ g \cdot (\alpha^{-1} m_\alpha)_{[\alpha]} = (g \cdot \alpha^{-1} m_\alpha)_{[g\alpha]}. \]
    The map is given by 
    \[ \phi \mapsto (\alpha^{-1} \phi(\alpha))_{[\alpha]}. \]
  \end{remark}

  \begin{lemma}
    \begin{enumerate}
      \item For all $G$-modules $M$ and $H$-modules $N$, we have a natural isomorphism
      \[ \Hom_{G}(M, \upoperatorname{Ind}_H^G N) \cong \Hom_{H}(M, N). \]
      \item The functor $\upoperatorname{Ind}_H^G: \Mod_H \to \Mod_G$ is exact.
    \end{enumerate}
  \end{lemma}
  \begin{proof}
    % \paragraph{(a)} Let $\alpha: M \to \upoperatorname{Ind}_H^G N$ be a homomorphism of $H$-modules.
    % We define $\Phi(\alpha): M \to N$ by $\Phi(\alpha)(m) := \alpha(m)(1)$ for all $m \in M$.
    % Conversely, let $\beta: M \to N$ be a homomorphism of $H$-modules.
    % We define $\Psi(\beta): M \to \upoperatorname{Ind}_H^G N$ by $\Psi(\beta)(m): g \mapsto \beta(g \cdot m)$ for all $m \in M$.
    % One can check $\Phi$ and $\Psi$ are well-defined and inverse to each other.

    % \paragraph{(b)} Let $0 \to M \to N \to P \to 0$ be an exact sequence of $H$-modules.
    % We need to show 
    % \[ 0 \to \upoperatorname{Ind}_H^G M \to \upoperatorname{Ind}_H^G N \to \upoperatorname{Ind}_H^G P \to 0 \] 
    % is exact.
    % By the \Yang{remark} and the exactness of product, we are done.
  \end{proof}

  Let $\phi: \upoperatorname{Ind}_H^G N \to N$ be the homomorphism of $H$-modules defined by $\phi(\phi') := \phi'(1)$ for all $\phi' \in \upoperatorname{Ind}_H^G N$.
  Then \Yang{lemma} gives a universal property of $\upoperatorname{Ind}_H^G N$: for all $G$-modules $M$ and homomorphisms $\alpha: M \to N$ of $H$-modules, there exists a unique homomorphism $\beta: M \to \upoperatorname{Ind}_H^G N$ of $G$-modules such that $\alpha = \phi \circ \beta$.
  \[ \begin{tikzcd}
    M \ar[r, "\beta", dashed] \ar[rd, "\alpha"'] & \upoperatorname{Ind}_H^G N \ar[d, "\phi"] \\
    & N
  \end{tikzcd} \]

  When $H = \{1\}$ is the trivial subgroup, we write $\upoperatorname{Ind}^G N$ for $\upoperatorname{Ind}_H^G N$.
  Thus $\upoperatorname{Ind}^G N = \Hom_{\Ab}(\bbZ[G], N)$.
  A $G$-module $M$ is called \emph{induced} if $M \cong \upoperatorname{Ind}^G N$ for some abelian group $N$.


  
